\chapter{Den tomme formelen}

\begin{motto}
Ein lov som ikkje kan vere usann, kan heller ikkje vere sann.\\ Eit heilt fag forveksla det fyrste med det andre, og kalla det eit fundament.
\end{motto}

Kvar moden disiplin har ein setning ho lærer fyrst, og som resten heng på. Fysikken har at kraft er masse gonger akselerasjon; biologien at form og funksjon møtest gjennom seleksjon. Designfaget har ein setning av same slag, og ho er eldre enn ordet «design» i si noverande tyding: \emph{form følgjer funksjon}. Ho er over hundre år gamal, ho vert tilskriven Louis Sullivan, ho vart slagordet for funksjonalismen og hjartet i modernismen, og ho er framleis det næraste faget kjem ein lov. Ho er òg, ved nærare ettersyn, ingen lov. Ho er ein tautologi forkledd som ein lov. Og oppdaginga av at ho er tom, gjord grundigast av Jan Michl\autocite{michl1995}, er den einaste store sjølvkritiske innsikta faget nokon gong har produsert om sitt eige fundament. At innsikta er korrekt, og at faget likevel ikkje bygde noko nytt på henne, fortel det meste om disiplinen sin evne til å lære av seg sjølv.

\section{Mannen, pulten og ørna}

Setninga vart ikkje meisla ut som ein lov; ho vart runa fram som ein hymne, og det forklarer mykje. Vi er i Chicago, vinteren 1896. Louis Sullivan, fyrti år gamal, er på høgda av karrieren, mannen som har lært skyskraparen å vere vakker og kledd dei høge kontorbygningane i vertikale band som dreg auget oppover. Han sit og skriv ein artikkel for \emph{Lippincott's Magazine} om nett dette: korleis ein skal forme det høge kontorbygget kunstnarleg\autocite{sullivan1896tall}. Og midt i denne praktiske, nesten yrkestekniske drøftinga byrjar språket plutseleg å løfte seg. Sullivan held opp med å snakke om kontor og heisar og byrjar å snakke om naturen.

Sjå på sjølve passasjen, for han er avslørande. «Det er den allgjennomtrengjande lova for alt organisk og uorganisk, for alt fysisk og metafysisk, for alt menneskeleg og alt overmenneskeleg,» skriv han, «for alle sanne ytringar av hovud, hjarte og sjel, at livet kjenner seg att i uttrykket, at form ever follows function. Dette er lova.» Og kva døme grip han til? Ikkje pulten, ikkje heisen, ikkje den ståltrukne berekonstruksjonen han faktisk arbeidde med. Han grip til ørna i flukt, til epleblomen på greina, til den arbeidande hesten, den lystige svana, den forgreina eika, det bukta vatnet ved kjeldevæld, til skyene og den rennande sola. Dette er ikkje arkitektur. Det er Walt Whitman omsett til arkitekturteori; det er romantisk naturmystikk der bygningen, knapt nemnd, skal vere like sjølvsagt uttrykk for sitt føremål som ørna er uttrykk for flukta. «Where function does not change, form does not change,» legg han til, som om eit kontorbygg var ein art i naturen, underlagt same naudsynet som fuglen sin venge.

For Sullivan var dette ein naturlov, og han trudde ikkje han formulerte ein etterprøvbar regel om korleis ein skulle teikne. Han formulerte ein tru: at det rette og det fagre fell saman i naturen, og at bygningen skal slutte seg til den same harmonien. Ein kan lese heile passasjen og ikkje finne éi einaste føresegn ein arkitekt kunne følgje, eller eit einaste tilfelle setninga ville utelukke. Men setninga hadde ein lagnad Sullivan ikkje kunne sjå: ho var kort, ho lét seg hugse, og ho kunne lausrivast frå organismen sin og gjerast til ein doktrine. Det var i denne tilstanden, som hymne og ikkje som hypotese, at frasen vart send vidare til eit heilt århundre som skulle ta henne for ein vitskapleg grunnsetning.

\section{Maskina å bu i, og villaen som var smak}

Tjuesju år seinare, i Paris, gjorde Le Corbusier det Sullivan aldri hadde tenkt: han hardna metaforen frå organisk til mekanisk og gjorde han til kampprogram. I \emph{Vers une architecture} frå 1923 fell den setninga som skulle gå verda rundt: huset er «une machine à habiter», ei maskin å bu i\autocite{lecorbusier1923}. Boka er full av fotografi, ikkje av hus, men av oceandampskip, av jagarfly, av Delage-bilen, sette ved sida av Parthenon, som om dei alle høyrde til same familie av rein, føremålstent form. Bodskapen er klar: slik flyet er reinsa for alt som ikkje tener flukta, skal huset reinsast for alt som ikkje tener bualivet. Bort med gesimsen, bort med taklista, bort med det skråe taket; fram med den nakne, kvite, geometriske flata.

Men sjå nøye på dei døma han valde. Flyet, bilen og dampskipet er objekt der eit hardt, målbart krav, å lette frå bakken, å gå fort, å flyte, faktisk pressar forma i ei bestemt retning. Huset har ikkje noko slikt krav. Eit hus kan stå like godt med skråtak som med flatt, med ornament som utan, og verda er full av bustader av begge slag som har halde i hundreåra. Le Corbusier lånte autoriteten frå maskinar der funksjonen verkeleg bind forma, og overførte han til eit felt der han ikkje gjer det. Maskin-metaforen var ikkje ein observasjon; han var eit retorisk grep, og eit svært verknadsfullt eit.

Sjå så kva som hender når vi held programmet opp mot det Le Corbusier faktisk bygde. I 1931 stod Villa Savoye ferdig i Poissy utanfor Paris, den reinaste utkrystalliseringa av læra, ein kvit kasse heva på tynne søyler, med band av vindauge og eit flatt tak. Han vart straks det helagaste ikonet i moderne arkitektur, og vert framleis lært som funksjonalismens fullkomne uttrykk. Men dei som skulle \emph{bu} der, fortel ei anna soge. Det flate taket lak nesten frå dag éin; vatn rann ned i romma, fukten åt seg innover, og sonen i familien Savoye vart sjuk i det kalde, klamme huset. Brevvekslinga mellom byggherren og arkitekten vart til slutt så bitter at familien truga med søksmål. Det flate taket har ingen funksjonell fordel framfor det skråe i eit klima med regn og snø; tvert imot er det vanskelegare å tette og lettare å øydeleggje. Kvifor var det då flatt? Fordi den geometriske, horisontale, ornamentfrie reinleiken var det moderne biletet på framsteg og ærlegdom, ein smak, vakker og kraftfull, men ein smak. Den forma som mest av alt skulle vere reint svar på bualivets funksjon, var altså den forma som dårlegast huste folk.

Le Corbusiers eige språk røpa kva som eigentleg gjekk føre seg, om ein berre les det nøye. «Arkitektur er det kunnige, korrekte og storslåtte spelet av volum samla i lys,» skreiv han i same bok, ein definisjon som er reint estetisk, om volum og lys og spel, og ikkje eit ord om bruk, husvære eller funksjon. Med éi hand kunngjorde han at huset var ei maskin; med den andre at arkitektur var eit spel av former i lys. Begge utsegnene er hans, og dei lever side om side utan at motsetnaden vert spurd: ein retorikk om funksjon og ein praksis av smak, haldne saman av ein frase som var tom nok til å romme begge.

Same grepet hadde Adolf Loos gjort alt nokre år tidlegare, då han kring 1910 steig opp på podiet for å halde det føredraget som vart kjent som «Ornament und Verbrechen», ornament og brotsverk\autocite{loos1908ornament}. Utsmykking var teikn på ein primitiv kultur, slo han fast, og den moderne mannen som tatoverte seg var anten ein forbrytar eller ein degenerert. Ornamentet var ikkje berre overflødig; det var umoralsk, eit sløseri med arbeidskraft og materiale, eit spor av barbariet i den siviliserte verda. Men sjå kva slags argument dette eigentleg er. Loos grunngjev ikkje fråvêret av ornament med at det tener bruken betre, for eit glatt hus husar ikkje betre enn eit utsmykka. Han grunngjev det moralsk og kulturelt: ornament er tilbakeståande, reinleik er framskriden. Det er ein \emph{smaksdom} kledd i sedelærdsspråk. Frå Loos til Le Corbusier går den same tråden: ein smak som kallar seg ei naudsyn, og som retorikken om funksjon kledde i naudsynets språk.

\begin{kasus}{Villa Savoye, det lekkande taket og smaken som kalla seg funksjon}
Villa Savoye er Michls argument støypt i betong. Le Corbusier hadde lånt autoriteten frå flyet og dampskipet, der ein hard, målbar straffefunksjon verkeleg pressar forma i éi retning, og bore han over til villaen, der ingen slik straff finst, bortsett frå den straffa naturen sjølv delte ut då regnet kom og taket svikta. Når funksjonen verkeleg bind forma, ser vi det fordi forma fungerer; her batt han henne ikkje, forma kom frå smaken, og resultatet lak. At det fremste byggverket i heile læra om at form følgjer funksjon ikkje fungerte for dei som budde i det, er ikkje ei pinleg fotnote. Det er sjølve diagnosen: setninga skjulte at valet var eit estetisk val, og kalla det ein konsekvens av bruken, heilt til bruken sjølv, i form av ein sjuk gut og eit vått golv, gav setninga ly.
\end{kasus}

\section{Michls disjunksjon}

Her kjem innvendinga som skulle ha endra alt. Michl stilte eit enkelt spørsmål til formelen: kva tyder «funksjon»? Og han viste at berre to svar er moglege, og at begge gjer setninga tom\autocite{michl1995}.

Det fyrste svaret er at funksjon tyder \emph{det objektet skal gjere} i snever forstand: stolen skal bere ein sitjande kropp, kanna skal halde og helle væske. Men dette underbestemmer forma fullstendig. Tusenvis av radikalt ulike stolar ber den same kroppen like godt: pinnestolen, lenestolen, krakken, barokkstolen, røyrstolen, plaststolen i hagen. Dei løyser alle sittefunksjonen, og likevel deler dei knapt ei linje. Om forma verkeleg følgde funksjonen i denne tydinga, skulle alle stolar ha konvergert mot éi optimal form. Dei gjer det ikkje, og dei har aldri gjort det. Altså følgjer forma ikkje funksjonen i snever forstand; funksjonen seier alt for lite til å bestemme henne.

Det andre svaret reddar formelen ved å utvide «funksjon» til å tyde alt forma gjer: ikkje berre sittinga, men det kulturelle uttrykket, statusen, den estetiske gleda, produksjonsøkonomien, tidsånda. No stemmer setninga alltid, fordi kva som helst ei form gjer kan kallast funksjonen ho følgjer. Barokkstolen følgjer sin funksjon, som er å syne makt; plaststolen følgjer sin, som er å vere billeg; designarstolen følgjer sin, som er å signalisere smak. Men ein påstand som er sann uansett kva du observerer, fortel deg ingenting før du observerer. Han kan ikkje føreseie ei einaste form, ikkje skilje ein god form frå ein dårleg, ikkje rettleie eit einaste val. Han er ikkje gal; han er tom. Han har forma til ei forklaring og innhaldet til ein sirkel.

Mellom desse to lesingane finst det inga tredje som både er ikkje-triviell og sann. Anten tyder funksjon noko presist, og då er formelen \emph{usann}, fordi forma er underbestemt; eller funksjon tyder alt, og då er formelen \emph{sann men tom}. Dette er ikkje retorikk. Det er den logiske strukturen til faget sin grunnsetning, og strukturen er den til ein tautologi. Michl hadde alt på syttitalet og åttitalet kjent på dette i artiklane sine om postmodernismen\autocite{michl1989}, og han kom tilbake til det gong på gong, til sist i ein hard dom over heile det modernistiske regimet i designutdanninga\autocite{michl2014}.

Det er verdt å dvele litt ved kor reint logisk poenget er, for det vert ofte mildna i refereringa. Folk les Michl som om han sa at funksjonalismen overdreiv, eller at funksjon er éin faktor blant fleire som verkar på forma. Men det er ikkje det han viser. Han viser noko strengare: at formelen, slik han faktisk vert brukt, ikkje kan vere informativ, av reint språkleg grunn. Ein setning gjev informasjon berre i den grad ho utelukkar noko, i den grad det finst tilstandar i verda ho seier ikkje kan vere. «Det regnar» gjev informasjon fordi det utelukkar tørt vêr. Spør no kva «form følgjer funksjon» utelukkar. I den vide lesinga: ingenting. Det finst inga tenkjeleg form som ikkje kan seiast å følgje sin funksjon, så snart funksjon får tyde alt forma gjer. Det er dette som er forskjellen mellom ein lov og ein tautologi, og det er denne forskjellen faget aldri tok inn over seg. Eit heilt århundre med formgjevarar vart oppdregne på ein regel som ikkje kunne brytast, og forveksla det med ein regel som ikkje kunne feile.

\section{Kva tomleiken skjulte}

Det verste er ikkje at formelen var tom. Det verste er kva tomleiken løynde. Når du seier at forma følgjer funksjonen, og funksjon tyder alt, har du i røynda sagt at forma følgjer \emph{noko}, og latt vere usagt kva, kven som avgjer det, og kva som vert lagt til side. Formelen gav modernismen eit alibi for å presentere reine estetiske og ideologiske val som naudsynte konsekvensar av funksjonen. Det flate taket, den kvite veggen, det blanke fråvêret av ornament; alt vart framstilt som det funksjonen kravde, når det i røynda var det ein bestemt smak føretrekte. Funksjonalismen var ikkje funksjonell. Han var ein stil som hevda å vere fråvêret av stil, og fordi grunnsetninga kalla stilen «funksjon», kunne faget aldri drøfte han som det han var: eit val mellom andre moglege val, gjort på vegne av nokon, med fylgjer for andre.

Her ligg strukturen som går att gjennom heile denne boka, no synleg i sjølve hjartesetninga. Eit fundament skal gjere vala synlege og diskuterbare. Denne setninga gjer det motsette: ho gøymer vala bak ein retorikk om naudsyn. «Forma følgjer funksjonen» tyder i praksis «mine val er ikkje val, dei er det saka krev». Det er ikkje eit fundament. Det er eit slør, og det er vove så fint at den som ber det, sjølv trur han går naken og ærleg.

Opprøret kom innanfrå, og det vart òg ufarleggjort. Då Robert Venturi i 1966 svara modernismen med at han ville ha både–og framfor anten–eller, kompleksitet og motsetnad framfor renleik, var dette i sak ei stadfesting av Michls poeng frå den andre kanten\autocite{venturi1966}. Venturi sa at bygningen alltid gjer fleire ting på ein gong, at meininga er rik og motstridande, at det reine funksjonsspråket var ei forarming. Men faget tok ikkje dette som ei oppmoding til å byggje ein betre teori om kva forma faktisk svarar på. Det tok det som lisens til å sleppe kravet om grunngjeving heilt: om alt er kompleksitet og motsetnad, treng eg ikkje seie kvifor noko er betre enn noko anna. Modernismens falske naudsyn vart bytt ut med postmodernismens fritak frå grunngjeving. Ingen av delane var eit fundament; begge var måtar å sleppe unna eitt.

\section{«Less is more» og paviljongen utan funksjon}

Det finst ei setning som står ved sida av Sullivan si som modernismens andre store grunnformel, og ho er endå reinare i si tomleik, fordi ho ikkje eingong later som ho handlar om funksjon. «Less is more», sa Ludwig Mies van der Rohe, og gjorde reduksjonen sjølv til lov. Mindre er meir: jo færre element, jo reinare verknad; ta bort til det ikkje er meir å ta bort, og det som står att, ber ein verdigheit dei rikare formene ikkje når. Det er ein vakker tanke, og som dei fleste vakre tankar i dette faget er han umogleg å motseie, fordi han er umogleg å prøve. Mindre er meir, av kva? Meir verdig? Meir sant? Meir funksjonelt? Mies sa det aldri, fordi setninga verkar nett ved å la det stå ope. Ho er ikkje ein observasjon om at reduksjon tener noko bestemt føremål; ho er ei estetisk truvedkjenning kledd i forma til ein logisk lov, og ho deler den same tomleiken Michl avslørte i Sullivan si setning: ho utelukkar ingenting, og ho ber difor ingen informasjon.

Den klaraste utkrystalliseringa av denne læra reiste Mies i 1929, ikkje som eit hus, men som ein gest. Tyskland skulle ha ein nasjonal paviljong på verdsutstillinga i Barcelona, og Mies fekk oppdraget. Det han laga, var eit av dei mest reproduserte byggverka i heile det tjuande hundreåret: ein låg, horisontal plattform av travertinmarmor, frittståande veggar av grøn marmor og gyllen onyks som ikkje møttest i hjørne men gleid forbi kvarandre, eit tak som såg ut til å sveve på tynne, krossforma stålsøyler, eit stille vassspegel med ein einsleg skulptur. Og her er det avslørande: paviljongen hadde inga oppgåve. Han skulle ikkje huse noko, ikkje stille ut noko, ikkje verne nokon mot vêret. Han var bygd for å bli opna seremonielt av den spanske kongen, og elles for å vere eit reint biletet av den tyske kulturen, ein arkitektur utan annan funksjon enn å vere arkitektur. Det fanst ingen kraft som valde formene hans, ikkje eit regn å halde ute, ikkje ein kropp å huse. Kvar einaste line var fri smak, og resultatet var av ein slik skjønnheit at det vart sjølve definisjonen på det moderne rommet.

Funksjonalismen sin mest dyrka helgedom var altså ein bygning utan funksjon. Mies sjølv var ærleg nok til aldri å hevde at han søkte funksjon; han søkte det han kalla «beinahe nichts», nesten ingenting, ei nesten tom form av ytste presisjon. Men det er nett dette som røper kva heile rørsla dreiv med. Når den arkitekten som meir enn nokon legemleggjorde modernismens estetiske ideal, byggjer sitt meisterverk heilt utan bruk, og det likevel vert hylla som funksjonalismens høgdepunkt, då har faget sagt høgt at «funksjon» aldri var det det handla om. Det handla om ei bestemt form, open, horisontal, rein, dyr, og om eit ord som kunne kle den forma i naudsynets språk. Barcelona-paviljongen er Michl sitt argument bygd i marmor.

\section{Då ei rørsle vart ein etikett}

Korleis vart ein europeisk strid om kva forma skulle tene, til ein amerikansk stil ein kunne kjøpe? I 1932 opna det unge Museum of Modern Art i New York ei utstilling som skulle bli ei av dei mest innflytelsesrike i designhistoria. Kuratorane var ein ung arkitekturhistorikar, Henry-Russell Hitchcock, og ein endå yngre, velståande estet ved namn Philip Johnson. Dei hadde reist gjennom Europa og sett den nye arkitekturen, Gropius, Mies, Le Corbusier, Oud, og dei kom heim med ein tese: at alle desse spreidde, stridande, ofte politisk radikale europeiske eksperimenta eigentleg var greiner av éin ny stil. Den stilen gav dei eit namn som vart hengjande ved: \emph{den internasjonale stilen}\autocite{hitchcock1932}.

Legg merke til kva som faktisk skjedde i denne omdøypinga. Dei europeiske modernistane hadde, dei fleste av dei, nekta at dei dreiv med stil i det heile. Stil var nett det dei meinte å ha overvunne; dei bygde, hevda dei, ut frå funksjon, ut frå sosial naud, ut frå materialet og produksjonen, ikkje ut frå smak. For mange av dei, og særleg for Hannes Meyer-fløya, var arkitekturen ein politisk handling, eit svar på arbeidarklassen sin bustadsnaud. Hitchcock og Johnson skar bort heile dette grunnlaget. Dei var ikkje interesserte i sosialismen, i funksjonen, i programmet; dei var interesserte i korleis bygningane såg ut. Og dei reduserte heile rørsla til tre formale kjenneteikn ein kunne katalogisere som ein botanikar katalogiserer artar: arkitektur som volum snarare enn masse, regelmessigheit snarare enn aksial symmetri, og fråvær av påsett ornament. Det var ei reint estetisk lesing, og ho var på sett og vis ærlegare enn modernistane sin eigen retorikk, for ho sa høgt at dette var ein måte ting såg ut på, ikkje ei naudsyn. Men ho gjorde det ved å tømme rørsla for alt innhald utanom utsjånaden.

Her ligg den doble ironien. For det fyrste: at det var \emph{mogleg} å destillere funksjonalismen til ein rein utsjånad utan funksjon, prova at funksjonen aldri hadde vore det berande, at det heile tida var ein form ein kunne skilje frå si påståtte grunngjeving og selje for seg. For det andre: Philip Johnson, mannen som gjorde funksjonalismen om til ein stiletikett, hadde sjølv ingen tru på funksjonen; han var ein estet som elska forma og sidan, utan blygsel, vart faget sin mest profilerte postmodernist då moten snudde. Den same mannen kunne selje funksjonalismen som stil og postmodernismen som stil, fordi for han var begge alltid berre det. At det var han, og ikkje ein motstandar, som sette namnet på rørsla, er diagnosen sjølv: funksjonalismen vart fyrst forstått til botnar av ein mann som ikkje trudde på funksjon, og som difor såg klart at det handla om utsjånad. Han kalla det den internasjonale stilen, og han hadde rett. Det var ein stil. Det var berre ikkje det rørsla hadde sagt om seg sjølv.

Den kvite boksen vart etter dette eit dogme. Det som hjå Le Corbusier var eit kampprogram mot ein heil borgarleg byggeskikk, vart hjå Hitchcock og Johnson eit sett formale reglar, og vart hjå dei tusen arkitektane som kom etter, ein vane: flatt tak, kvit vegg, band av vindauge, ingen pynt, gjort ikkje fordi funksjonen kravde det, men fordi det var slik moderne arkitektur såg ut, og fordi grunnsetninga om form og funksjon stod klar til å kle vana i naudsynets språk når nokon spurde. Reyner Banham såg dette tydeleg då han granska heile rørsla i ettertid: det modernismen kalla funksjonell logikk, var oftast ein formalistisk smak som hadde lært å snakke ingeniørens språk\autocite{banham1960}. Den kvite boksen var ikkje konklusjonen på ein vitskap om form. Han var ein stil som hadde gløymt at han var ein stil, og som difor kunne påleggjast som lov.

\section{Setninga var aldri Sullivans eige}

Det er ein lærdom gøymd alt i opphavssoga, og han gjer tomleiken endå pinlegare. For formelen var ikkje Sullivans oppfinning. Han var arvegods, og då Sullivan skreiv han ned i 1896, var tanken alt over hundre år gamal. Edward Robert De Zurko synte dette i 1957, i ei bok som faget burde ha lese og ikkje las: \emph{Origins of Functionalist Theory} sporar funksjonalistiske argument attende gjennom det nittande hundreåret, gjennom opplysingstida, heilt til antikken og den romerske arkitekturteorien\autocite{dezurko1957origins}. Tanken om at forma skal svare til føremålet er ikkje ein moderne innsikt; han er ein topos, ein retorisk fellesstad som har vore med oss sidan Cicero skreiv at det som er nyttig òg er fagert. Sullivan gjorde ikkje anna enn å gje den eldgamle frasen ein amerikansk, organisk klang.

Den næraste kjelda hans var ein landsmann to generasjonar eldre. Billedhoggaren Horatio Greenough, som døydde i 1852, skreiv på 1840-talet ei rekkje essay om form i naturen og kunsten der han nett hevda at skjønnheit er «løftet om funksjon» og at naturen aldri pyntar utan grunn. Greenough sitt yndlingsdøme var det amerikanske klipperskipet, og det er verdt å sjå kvifor, for det avslører heile gliden i eitt bilete. Klipperen er vakker, sa Greenough, fordi kvar line i skroget er tvinga fram av kravet om å skjere fort gjennom sjøen; ingenting er lagt til for syns skuld. Her har funksjonen verkeleg forma forma, fordi havet straffar kvar feil med tap av fart. Men sjå kva som hender når argumentet vert lyfta frå skipet til huset. Havet finst ikkje lenger. Det finst ingen kraft som straffar det skråe taket eller det utsmykka rekkverket med målbart tap, slik sjøen straffar det klumpete skroget. Greenough, og etter han Sullivan, og etter han Le Corbusier, tok den tvingande logikken frå eit objekt der naturen verkeleg vel, og bar han over til objekt der berre smaken vel. Klipperskipet er ikkje eit prov på at form følgjer funksjon. Det er eit prov på at \emph{når} det finst ein hard, målbar straffefunksjon, då, og berre då, vert forma bestemt, nett det Michl skulle formulere strengt halvannan hundreår seinare.

Ralph Waldo Emerson tok tanken vidare, og det er gjennom denne transcendentalistiske lina, Greenough, Emerson, Whitman, at Sullivan arva henne. Poenget er ikkje akademisk pedanteri om kven som sa kva fyrst. Poenget er kva det fortel at faget bygde sitt einaste «fundament» på ein frase det ikkje visste opphavet til, og som ved nærare ettersyn aldri var ein observasjon om bygningar i det heile, men ein gamal moralsk-estetisk lengt etter at det rette og det vakre skulle falle saman. Eit fag som ikkje kjenner soga til sin eigen grunnsetning, kan heller ikkje vite kva han eigentleg påstår, og kan difor ta ein arva smak for ei nyvunnen lov.

\section{Hundre år utan oppfølging}

Michls kritikk er ikkje obskur. Han er trykt, omsett, antologisert, undervist. Michl publiserte argumentet fyrst i 1989, i tidsskriftet \emph{Pro Forma}, og skjerpa det i 1995 til den kanoniske artikkelen «Form Follows WHAT?», seinare gjort fritt tilgjengeleg på nettet, der det har vorte ein av dei mest lesne enkelttekstane i heile fagkritikken. Tjuefem år med same argument, stadig klårare formulert, stadig breiare lese. Ingen har tilbakevist det. Det finst, så vidt eg veit, ikkje ein einaste publisert freistnad på å vise at Michl tek logisk feil, av den enkle grunn at han ikkje gjer det.

Og likevel: kva endra seg? Sjå i innføringsbøkene, sjå i pensumlistene. Sullivan-sitatet står framleis på fyrste sida, presentert utan åtvaring, som om det var ein etablert sanning og ikkje ein gjennomskoda tautologi. Michl, når han i det heile vert nemnd, dukkar opp som ei «kritisk røyst» i ein boks ved sida av, ei meining mellom meiningar, ikkje ei avgjerd. Dette er den avslørande forma resepsjonen tok: ikkje motstand, men \emph{absorpsjon}. Faget tilbakeviste ikkje Michl, for det kunne det ikkje. Det gjorde noko meir effektivt: det gjorde han til pensum. Ein kritikk som vert pensum utan å endre praksis, er uskadeleggjort, ikkje teken på alvor. Han er vaksinert inn i kroppen i ein dose for liten til å verke. Studenten lærer formelen på måndag og kritikken av formelen på tysdag, og ingen lærar ber han nokon gong gjere noko med motsetnaden, fordi det ville krevje eit svar faget ikkje har.

Det er instruktivt å samanlikne med korleis eit fag med fundament handsamar ein tilsvarande innvending. Då matematikaren Gottlob Frege fekk Bertrand Russells brev i 1902 som synte at grunnlaget hans for aritmetikken inneheldt ein motseiing, skreiv Frege at «aritmetikken vaklar» og braut av arbeidet sitt; heile fagfeltet kasta seg over problemet, og ut av krisa voks ny grunnlagsforsking som forma matematikken i hundreåret etter. Innvendinga var dødeleg, og faget handsama henne som dødeleg. Då Michl synte at designfagets grunnsetning er ein tautologi, vakla ingenting. Faget skreiv henne inn i ein lærebokboks og heldt fram. Skilnaden er ikkje at Michls argument er svakare enn Russells; strukturelt er dei i slekt, begge syner at eit påstått fundament er innhaldslaust eller motstridande. Skilnaden er at det eine faget hadde noko å tape, og difor reagerte, medan det andre ikkje hadde det, og difor ikkje reagerte.

For å skifte ut ein tom formel treng du ein som ikkje er tom. Du må kunne seie, presist, kva forma faktisk følgjer: kva som verkar på henne, korleis, og kvar dei verkande kreftene kjem frå. Det krev eit omgrep om forma som ein posisjon i eit rom av moglege former, om kreftene som mange samtidige trykk i det rommet, og om formgjevaren som ein som navigerer det. Michl hadde den negative innsikta, at den gamle formelen var tom, men ikkje den positive: kva ein skulle setje i staden. Og faget som helse hadde korkje det eine eller det andre. Det kunne diagnostisere tomleiken, og heldt så fram med å bruke det det visste var tomt, fordi det ikkje hadde noko anna.

Ein kan spørje kvifor formelen overlever, når han er så tydeleg gjennomskoda. Svaret ligg i kva han gjev, og det er ikkje sanning, men tryggleik. Den tomme formelen er bra å ha nett fordi han er tom: han kan rettferdiggjere kva som helst arbeid i ettertid, og han kan aldri ta deg i å ha gjort feil. Påstår nokon at arbeidet ditt er stygt eller skadeleg, kan du svare at det følgjer funksjonen, og sidan ingen kan seie kva funksjonen ikkje omfattar, kan ingen motseie deg. Ein setning med denne eigenskapen vert ikkje gjeven opp frivillig, for han løyser eit reelt problem for utøvaren, problemet med å måtte forsvare vala sine, sjølv om han løyser det ved å gjere forsvaret innhaldslaust. Eit fag som vernar om retten til aldri å kunne ta feil, har bytt bort retten til nokon gong å ha rett.

\laeringsmal{\begin{itemize}
\item Gjere greie for Michls disjunksjon og vise, med eigne døme, kvifor begge lesingane av «funksjon» gjer formelen anten usann eller tom.
\item Forklare skiljet mellom ein lov og ein tautologi ved hjelp av omgrepet informasjon som det å utelukke moglege tilstandar.
\item Spore funksjonalismens slagord attende før Sullivan, til Greenough og den eldre funksjonalistiske topos, og vurdere kva opphavet seier om kva frasen eigentleg påstår.
\item Skildre korleis faget tok imot Michls kritikk, og drøfte skilnaden mellom å tilbakevise, å absorbere og å ignorere ein grunnlagskritikk.
\end{itemize}}

\ovingar{\begin{enumerate}
\item Finn tre radikalt ulike gjenstandar som løyser same funksjon (til dømes tre måtar å sitje på, eller tre måtar å lyse opp eit rom). Vis korleis kvar av dei «følgjer funksjonen» like godt, og bruk dette til å demonstrere underbestemminga konkret.
\item Vel eit produkt og skriv to grunngjevingar for forma: éi som hevdar at ho følgjer funksjonen, og éi som røper det estetiske eller ideologiske valet bak. Drøft kva den fyrste grunngjevinga skjuler.
\item Samanlikn Greenough sitt klipperskip med eit moderne objekt utan tvingande fysisk straffefunksjon (ein stol, ei lampe, eit grensesnitt). Kvar finst «havet» som vel forma, og kvar finst det ikkje?
\item Diskuter i seminar: Om Michls kritikk er korrekt og uimotsagd, kvifor står Sullivan-sitatet framleis i innføringsbøkene? Kva ville det krevje av faget å fjerne det?
\end{enumerate}}

\vignett

Formelen var den intellektuelle kjernen i designutdanninga, og kjernen var hol. Men ein kunne svare at den verkelege ambisjonen aldri budde i slagord; at faget visste dette sjølv, og at det ein gong sette seg føre å erstatte slagordet med noko strengare, ein metode, ei vitskap, eit verkeleg fundament. Det forsøket fanst. Det hadde ei adresse, ein liten by i Sør-Tyskland, og det kom nærare enn nokon før eller sidan. Lat oss reise til Ulm.

\vidarelesing{\fullcite{sullivan1896tall}; \fullcite{lecorbusier1923}; \fullcite{loos1908ornament}; \fullcite{hitchcock1932}; \fullcite{banham1960}; \fullcite{michl1995}; \fullcite{venturi1966}}
